% Statement of Significance - Journal of Biomedical Informatics
% Place at the end of the Introduction section in main_jbi.tex
\documentclass[12pt]{article}
\usepackage[margin=1in]{geometry}
\usepackage{booktabs}
\usepackage{array}

\begin{document}

\section*{Statement of Significance}
\noindent
The authors summarize in 150 words or less the contribution of the paper to the existing literature.

\begin{table}[h]
\centering
\begin{tabular}{@{}p{0.28\textwidth}p{0.34\textwidth}p{0.34\textwidth}@{}}
\toprule
\textbf{Problem or Issue} & \textbf{What is Already Known} & \textbf{What this Paper Adds} \\
\midrule
Session-based EHR retrieval suffers from stale-context pollution after abrupt topic pivots during chart review. & Session-based retrieval degrades under context drift; JEPA-style predictive learning enables latent forecasting. & CMA combines SPD-manifold intent encoding, geodesic-shift gating, and JEPA prefetching to reduce clinical retrieval latency by 23.8\% while exposing the limits of sparse retrieval (TF-IDF and BM25). \\
\bottomrule
\end{tabular}
\end{table}

\end{document}
